\part{IV gimnazijos klasė}

\section{Integralas ir pirmykštė funkcija}\label{sec:integralas-ir-pirmykste-funkcija}

\begin{tcolorbox}[title=Primename iš ankstesnių klasių, colback=yellow!10]
Šiai temai reikia mokėti: \textbf{išvestinė ir diferencijavimas}, nagrinėta 11 klasėje.
Jei reikia, pasikartok: \ref{sec:isvestine-ir-diferencijavimas}.
\end{tcolorbox}

\subsection{Apibrėžtys}
\begin{apibreztis}
Funkcija \(F\) vadinama funkcijos \(f\) pirmykšte funkcija intervale \(I\), jeigu kiekvienam \(x\in I\) galioja \(F'(x)=f(x)\).
\end{apibreztis}
\begin{apibreztis}
Neapibrėžtinis integralas yra visų funkcijos \(f\) pirmykščių funkcijų šeima:
\[
\int f(x)\,dx=F(x)+C,\qquad C\in\mathbb{R}.
\]
Ženklas \(dx\) primena, pagal kurį kintamąjį integruojama.
\end{apibreztis}
\begin{apibreztis}
Pagrindinė integravimo užduotis yra atvirkštinė diferencijavimui: žinome kitimo greitį \(f(x)\), o ieškome sukaupto dydžio \(F(x)\).
\end{apibreztis}

\subsection{Teoremos ir savybės}
\begin{teorema}
Jeigu \(F\) yra viena funkcijos \(f\) pirmykštė intervale \(I\), tai visos \(f\) pirmykštės tame intervale yra \(F+C\), kur \(C\) yra konstanta.
\end{teorema}
\begin{proof}
Jei \(G'=f\) ir \(F'=f\), tai \((G-F)'=G'-F'=0\). Funkcija, kurios išvestinė intervale lygi nuliui, tame intervale yra pastovi, todėl \(G-F=C\), t. y. \(G=F+C\).
\end{proof}
\begin{savybe}
Integralas yra tiesinis:
\[
\int (a f(x)+b g(x))\,dx=a\int f(x)\,dx+b\int g(x)\,dx.
\]
\end{savybe}
\begin{savybe}
Dažniausios pirmykščių formulės:
\[
\int x^n\,dx=\frac{x^{n+1}}{n+1}+C\quad(n\ne -1),\qquad
\int \frac{1}{x}\,dx=\ln|x|+C,
\]
\[
\int e^x\,dx=e^x+C,\quad
\int \sin x\,dx=-\cos x+C,\quad
\int \cos x\,dx=\sin x+C.
\]
\end{savybe}
\begin{pastaba}
Išplėstiniame kurse dažnai reikia integruoti ir paprastas sudėtines funkcijas. Jeigu vidinės funkcijos išvestinė yra šalia, tinka keitimo mintis: \(\int f'(x)\varphi(f(x))\,dx=\Phi(f(x))+C\), kai \(\Phi'=\varphi\).
\end{pastaba}

\subsection{Taisyklės ir algoritmai}
\begin{enumerate}[label=\arabic*.]
\item Pertvarkykite reiškinį taip, kad matytųsi pagrindinės formulės: išskleiskite skliaustus, atskirkite trupmenas, laipsnius rašykite \(x^n\) pavidalu.
\item Integruokite kiekvieną dėmenį atskirai ir nepamirškite konstantų daugiklių.
\item Pridėkite \(+C\). Be šios konstantos atsakymas neaprašo visų pirmykščių.
\item Patikrinkite diferencijuodami gautą rezultatą. Jeigu išvestinė grįžta į pradinę funkciją, integralas rastas teisingai.
\end{enumerate}

\subsection{Iliustruojantys pavyzdžiai}
\begin{enumerate}[label=\arabic*.]
\item Raskime \(\int (6x^2-4x+5)\,dx\).
\begin{align}
\int (6x^2-4x+5)\,dx
&=6\cdot\frac{x^3}{3}-4\cdot\frac{x^2}{2}+5x+C\\
&=2x^3-2x^2+5x+C.
\end{align}
\textbf{Atsakymas:} \(2x^3-2x^2+5x+C\).

\item Raskime pirmykštę \(F\), jei \(F'(x)=3\sqrt{x}-\frac{2}{x}\), \(x>0\), ir \(F(1)=4\).
\begin{align}
F(x)&=\int \left(3x^{1/2}-2x^{-1}\right)\,dx\\
&=3\cdot\frac{x^{3/2}}{3/2}-2\ln x+C\\
&=2x^{3/2}-2\ln x+C.
\end{align}
Iš sąlygos \(F(1)=4\):
\[
2\cdot1^{3/2}-2\ln1+C=4,\qquad C=2.
\]
\textbf{Atsakymas:} \(F(x)=2x^{3/2}-2\ln x+2\), kai \(x>0\).

\item Apskaičiuokime \(\int 2x\cos(x^2+1)\,dx\).
\begin{align}
u&=x^2+1,\qquad du=2x\,dx,\\
\int 2x\cos(x^2+1)\,dx&=\int \cos u\,du=\sin u+C.
\end{align}
\textbf{Atsakymas:} \(\sin(x^2+1)+C\).
\end{enumerate}

\subsection{Dažnos klaidos}
\begin{itemize}
\item Pamirštama konstanta \(C\). Pavyzdžiui, \(\int 2x\,dx=x^2+C\), ne vien \(x^2\).
\item Taikoma laipsnio formulė atvejui \(n=-1\). \(\int x^{-1}\,dx\) yra \(\ln|x|+C\), o ne \(\frac{x^0}{0}\).
\item Sumai integruoti pradedama taikyti sandaugos taisyklės analogija. Integralas tiesinis sumai, bet \(\int f(x)g(x)\,dx\) paprastai nelygu \(\int f(x)\,dx\int g(x)\,dx\).
\end{itemize}

\subsection{Mini praktika}
\begin{enumerate}[label=\alph*)]
\item Raskite \(\int (4x^3-6x+2)\,dx\).
\item Raskite \(F\), jei \(F'(x)=\frac{1}{x}+e^x\), \(x>0\), ir \(F(1)=e\).
\item Patikrinkite diferencijuodami: ar \(-3\cos x+2e^x+C\) yra funkcijos \(3\sin x+2e^x\) pirmykštė?
\end{enumerate}

\subsection{Ryšys su kitomis temomis}
Pirmykštė funkcija yra tiltas nuo išvestinės prie apibrėžtinio integralo. Jei išvestinė aprašo momentinį kitimą, integralas aprašo sukauptą pokytį, plotą, tūrį ar bendrą darbą.

\section{Integralo taikymai}\label{sec:integralo-taikymai}

\subsection{Apibrėžtys}
\begin{apibreztis}
Apibrėžtinis integralas \(\int_a^b f(x)\,dx\) yra integralinių sumų riba. Geometriškai, kai \(f(x)\ge0\), jis lygus kreivinės trapecijos plotui po grafiku intervale \([a;b]\).
\end{apibreztis}
\begin{apibreztis}
Orientuotas plotas reiškia, kad grafiko dalys virš \(x\) ašies duoda teigiamą indėlį, o dalys žemiau \(x\) ašies -- neigiamą.
\end{apibreztis}
\begin{apibreztis}
Funkcijos vidutinė reikšmė intervale \([a;b]\) yra
\[
\overline f=\frac{1}{b-a}\int_a^b f(x)\,dx.
\]
\end{apibreztis}

\subsection{Teoremos ir savybės}
\begin{teorema}[Niutono ir Leibnico formulė]
Jeigu \(F'=f\) intervale \([a;b]\), tai
\[
\int_a^b f(x)\,dx=F(b)-F(a).
\]
\end{teorema}
\begin{proof}
Pirmykštė funkcija \(F\) aprašo sukauptą dydį. Integralas sudeda mažus pokyčius \(f(x)\Delta x\), todėl visa suma artėja prie bendro pokyčio \(F(b)-F(a)\).
\end{proof}
\begin{savybe}
Jeigu \(a<c<b\), tai \(\int_a^b f(x)\,dx=\int_a^c f(x)\,dx+\int_c^b f(x)\,dx\).
\end{savybe}
\begin{savybe}
Jeigu \(f(x)\ge g(x)\) intervale \([a;b]\), tai plotas tarp grafikų yra
\[
S=\int_a^b (f(x)-g(x))\,dx.
\]
\end{savybe}
\begin{savybe}
Išplėstiniame kurse sukinio, gauto sukant \(y=f(x)\ge0\) apie \(x\) ašį, tūris:
\[
V=\pi\int_a^b (f(x))^2\,dx.
\]
\end{savybe}

\subsection{Taisyklės ir algoritmai}
\begin{enumerate}[label=\arabic*.]
\item Nustatykite, ką skaičiuojate: orientuotą integralą, tikrą plotą, vidutinę reikšmę ar tūrį.
\item Jei reikia ploto, raskite susikirtimo taškus ir nuspręskite, kuri funkcija viršuje.
\item Užrašykite integralą su teisingomis ribomis.
\item Raskite pirmykštę ir pritaikykite Niutono ir Leibnico formulę.
\item Patikrinkite ženklą ir vienetus. Plotas negali būti neigiamas, tūrio vienetai yra kubiniai.
\end{enumerate}

\subsection{Iliustruojantys pavyzdžiai}
\begin{enumerate}[label=\arabic*.]
\item Apskaičiuokime \(\int_0^3 (2x+1)\,dx\).
\begin{align}
\int_0^3 (2x+1)\,dx
&=\left[x^2+x\right]_0^3\\
&=(9+3)-0=12.
\end{align}
\textbf{Atsakymas:} \(12\).

\item Raskime plotą tarp \(y=x\) ir \(y=x^2\).
Susikirtimo taškai:
\[
x=x^2,\qquad x(x-1)=0,\qquad x=0,\;x=1.
\]
Kai \(0<x<1\), viršuje yra \(y=x\), todėl
\begin{align}
S&=\int_0^1 (x-x^2)\,dx\\
&=\left[\frac{x^2}{2}-\frac{x^3}{3}\right]_0^1
=\frac12-\frac13=\frac16.
\end{align}
\textbf{Atsakymas:} \(\frac16\).

\item Kreivė \(y=\sqrt{x}\) sukama apie \(x\) ašį, \(0\le x\le4\). Raskime sukinio tūrį.
\begin{align}
V&=\pi\int_0^4 (\sqrt{x})^2\,dx
=\pi\int_0^4 x\,dx\\
&=\pi\left[\frac{x^2}{2}\right]_0^4=8\pi.
\end{align}
\textbf{Atsakymas:} \(8\pi\).
\end{enumerate}

\subsection{Dažnos klaidos}
\begin{itemize}
\item Plotas tarp grafikų skaičiuojamas kaip \(\int (g-f)\,dx\), nors \(g\) yra apatinė funkcija. Tada gaunamas neigiamas skaičius.
\item Pamirštama perskirti intervalą, kai grafikas kerta \(x\) ašį. Tikras plotas yra atskirų teigiamų plotų suma.
\item Tūrio formulėje pamirštama kvadratu kelti funkciją: \(V=\pi\int (f(x))^2\,dx\), ne \(\pi\int f(x)\,dx\).
\end{itemize}

\subsection{Mini praktika}
\begin{enumerate}[label=\alph*)]
\item Apskaičiuokite \(\int_1^2 (3x^2-1)\,dx\).
\item Raskite plotą tarp \(y=4-x^2\) ir \(x\) ašies.
\item Raskite funkcijos \(f(x)=x^2\) vidutinę reikšmę intervale \([0;3]\).
\end{enumerate}

\subsection{Ryšys su kitomis temomis}
Integralo taikymai jungia funkcijų analizę, geometriją, fizikinius modelius ir tikimybinius tankius. Egzamino uždaviniuose dažnai reikia kartu rasti susikirtimo taškus, nubraižyti eskizą ir tik tada integruoti.

\section{Diferencialinės lygtys ir modeliai}\label{sec:diferencialines-lygtys-ir-modeliai}

\begin{tcolorbox}[title=Primename iš ankstesnių klasių, colback=yellow!10]
Šiai temai reikia mokėti: \textbf{išvestinės taikymai}, nagrinėta 11 klasėje.
Jei reikia, pasikartok: \ref{sec:isvestines-taikymai}.
\end{tcolorbox}

\subsection{Apibrėžtys}
\begin{apibreztis}
Diferencialinė lygtis yra lygtis, kurioje nežinoma funkcija siejama su jos išvestinėmis, pavyzdžiui, \(y'=2x\) arba \(N'=0{,}03N\).
\end{apibreztis}
\begin{apibreztis}
Bendrasis sprendinys aprašo visą sprendinių šeimą, o konkretusis sprendinys gaunamas pritaikius pradinę sąlygą, pavyzdžiui, \(y(0)=5\).
\end{apibreztis}

\subsection{Teoremos ir savybės}
\begin{teorema}
Lygties \(y'=ky\) sprendiniai yra \(y=Ce^{kx}\), kur \(C\) yra konstanta.
\end{teorema}
\begin{proof}
Jei \(y=Ce^{kx}\), tai \(y'=kCe^{kx}=ky\). Atvirkščiai, kai \(y\ne0\), galima atskirti kintamuosius:
\[
\frac{dy}{y}=k\,dx,\qquad \ln|y|=kx+C_1,\qquad y=Ce^{kx}.
\]
Sprendinys \(y=0\) taip pat įeina, kai \(C=0\).
\end{proof}
\begin{savybe}
Jeigu kitimo greitis proporcingas pačiam kiekiui, modelis yra eksponentinis: augimas, kai \(k>0\), ir nykimas, kai \(k<0\).
\end{savybe}
\begin{savybe}
Lygtis \(y'=f(x)\) sprendžiama integruojant: \(y=\int f(x)\,dx\).
\end{savybe}

\subsection{Taisyklės ir algoritmai}
\begin{enumerate}[label=\arabic*.]
\item Įvardykite nežinomą funkciją ir kintamąjį: \(N(t)\), \(s(t)\), \(T(t)\) ar pan.
\item Sąlygos žodžius perkelkite į lygtį: ``greitis proporcingas kiekiui'' reiškia \(N'=kN\).
\item Išspręskite lygtį integruodami arba atskirdami kintamuosius.
\item Pritaikykite pradinę sąlygą ir patikrinkite, ar atsakymas turi prasmę pagal kontekstą.
\end{enumerate}

\subsection{Iliustruojantys pavyzdžiai}
\begin{enumerate}[label=\arabic*.]
\item Išspręskime \(y'=4x-3\), \(y(1)=2\).
\begin{align}
y&=\int(4x-3)\,dx=2x^2-3x+C,\\
2&=2\cdot1^2-3\cdot1+C=-1+C,\\
C&=3.
\end{align}
\textbf{Atsakymas:} \(y=2x^2-3x+3\).

\item Populiacija auga pagal \(N'=0{,}08N\), \(N(0)=500\). Raskime \(N(t)\).
\[
N(t)=Ce^{0{,}08t},\qquad N(0)=C=500.
\]
\textbf{Atsakymas:} \(N(t)=500e^{0{,}08t}\).

\item Vaisto kiekis kraujyje mažėja pagal \(A'=-0{,}2A\), \(A(0)=40\). Kada liks \(10\) mg?
\begin{align}
A(t)&=40e^{-0{,}2t},\\
10&=40e^{-0{,}2t},\\
\frac14&=e^{-0{,}2t},\\
t&=\frac{\ln4}{0{,}2}\approx6{,}93.
\end{align}
\textbf{Atsakymas:} maždaug po \(6{,}93\) laiko vieneto.
\end{enumerate}

\subsection{Dažnos klaidos}
\begin{itemize}
\item Painiojama \(y'=k\) ir \(y'=ky\). Pirmuoju atveju gaunama tiesinė funkcija \(y=kx+C\), antruoju -- eksponentinė.
\item Pradinė sąlyga pritaikoma dar prieš randant bendrąjį sprendinį, todėl prarandama konstanta.
\item Eksponentinio nykimo modelyje pamirštamas neigiamas ženklas.
\end{itemize}

\subsection{Mini praktika}
\begin{enumerate}[label=\alph*)]
\item Išspręskite \(y'=6x^2\), \(y(0)=1\).
\item Išspręskite \(P'=0{,}05P\), \(P(0)=1200\).
\item Ką reiškia koeficientas \(k=-0{,}1\) modelyje \(Q'=kQ\)?
\end{enumerate}

\subsection{Ryšys su kitomis temomis}
Diferencialinės lygtys parodo, kaip išvestinė tampa modeliu. Jos remiasi integralais, eksponentinėmis ir logaritminėmis funkcijomis, o taikomuosiuose uždaviniuose dažnai susijungia su procentiniu kitimu.

\section{Kompleksiniai skaičiai}\label{sec:kompleksiniai-skaiciai}

\subsection{Apibrėžtys}
\begin{apibreztis}
Kompleksinis skaičius yra \(z=a+bi\), kur \(a,b\in\mathbb{R}\), o \(i^2=-1\). Skaičius \(a\) vadinamas realiąja dalimi, \(b\) -- menamąja dalimi.
\end{apibreztis}
\begin{apibreztis}
Skaičiaus \(z=a+bi\) modulis yra \(|z|=\sqrt{a^2+b^2}\), o jungtinis skaičius yra \(\overline z=a-bi\).
\end{apibreztis}
\begin{apibreztis}
Kompleksinėje plokštumoje skaičius \(a+bi\) vaizduojamas tašku \((a;b)\). Modulis yra atstumas nuo pradžios taško iki šio taško.
\end{apibreztis}

\subsection{Teoremos ir savybės}
\begin{savybe}
\[
z\overline z=|z|^2.
\]
\end{savybe}
\begin{proof}
\[
(a+bi)(a-bi)=a^2-b^2i^2=a^2+b^2=|z|^2.
\]
\end{proof}
\begin{teorema}
Kvadratinė lygtis \(ax^2+bx+c=0\), \(a\ne0\), su realiais koeficientais ir neigiamu diskriminantu turi dvi jungtines kompleksines šaknis:
\[
x=\frac{-b\pm i\sqrt{-D}}{2a},\qquad D=b^2-4ac<0.
\]
\end{teorema}
\begin{savybe}
Kompleksiniai skaičiai sudedami ir atimami dalimis, o dauginant naudojama taisyklė \(i^2=-1\).
\end{savybe}

\subsection{Taisyklės ir algoritmai}
\begin{enumerate}[label=\arabic*.]
\item Sudėdami ar atimdami grupuokite realiąsias ir menamąsias dalis.
\item Daugindami išskleiskite skliaustus ir kiekvieną \(i^2\) pakeiskite \(-1\).
\item Dalydami iš \(a+bi\), dauginkite skaitiklį ir vardiklį iš jungtinio \(a-bi\).
\item Spręsdami kvadratinę lygtį su \(D<0\), rašykite \(\sqrt{D}=i\sqrt{-D}\).
\end{enumerate}

\subsection{Iliustruojantys pavyzdžiai}
\begin{enumerate}[label=\arabic*.]
\item Apskaičiuokime \((2-3i)+(5+7i)\).
\[
(2-3i)+(5+7i)=7+4i.
\]
\textbf{Atsakymas:} \(7+4i\).

\item Padalykime \(\frac{3+i}{2-i}\).
\begin{align}
\frac{3+i}{2-i}
&=\frac{(3+i)(2+i)}{(2-i)(2+i)}\\
&=\frac{6+3i+2i+i^2}{4-i^2}
=\frac{5+5i}{5}=1+i.
\end{align}
\textbf{Atsakymas:} \(1+i\).

\item Išspręskime \(x^2-2x+5=0\).
\begin{align}
D&=(-2)^2-4\cdot1\cdot5=-16,\\
x&=\frac{2\pm i\sqrt{16}}{2}=1\pm2i.
\end{align}
\textbf{Atsakymas:} \(x=1-2i\) ir \(x=1+2i\).
\end{enumerate}

\subsection{Dažnos klaidos}
\begin{itemize}
\item Rašoma \(\sqrt{-9}=-3\). Teisingai: \(\sqrt{-9}=3i\) kompleksinių skaičių sistemoje.
\item Dalijant nepanaikinama menamoji dalis vardiklyje.
\item Pamirštama, kad \((bi)^2=b^2i^2=-b^2\).
\end{itemize}

\subsection{Mini praktika}
\begin{enumerate}[label=\alph*)]
\item Apskaičiuokite \((4+i)(1-3i)\).
\item Raskite \(|-6+8i|\).
\item Išspręskite \(x^2+6x+13=0\).
\end{enumerate}

\subsection{Ryšys su kitomis temomis}
Kompleksiniai skaičiai išplečia realiųjų skaičių aibę taip, kad kvadratinės lygtys visada turėtų šaknis. Jie ypač naudingi algebroje, trigonometrijoje ir vėlesnėse technologijų bei fizikos temose.

\section{Matricos ir tiesinės sistemos}\label{sec:matricos-ir-tiesines-sistemos}

\begin{tcolorbox}[title=Primename iš ankstesnių klasių, colback=yellow!10]
Šiai temai reikia mokėti: \textbf{lygčių ir nelygybių sistemos}, nagrinėta 9--10 klasėse.
Jei reikia, pasikartok: \ref{sec:lygciu-ir-nelygybiu-sistemos}.
\end{tcolorbox}

\subsection{Apibrėžtys}
\begin{apibreztis}
Matrica yra stačiakampė skaičių lentelė. Matricos elementas \(a_{ij}\) yra \(i\)-oje eilutėje ir \(j\)-ame stulpelyje.
\end{apibreztis}
\begin{apibreztis}
Tiesinė sistema yra lygčių sistema, kurioje visi nežinomieji įeina pirmuoju laipsniu.
\end{apibreztis}
\begin{apibreztis}
Determinantas \(2\times2\) matricai skaičiuojamas taip:
\[
\begin{vmatrix}a&b\\c&d\end{vmatrix}=ad-bc.
\]
\end{apibreztis}

\subsection{Teoremos ir savybės}
\begin{savybe}
Matricas galima sudėti tik tada, kai jų matmenys vienodi. Sudedami atitinkami elementai.
\end{savybe}
\begin{savybe}
Matricų sandauga \(AB\) apibrėžta tada, kai \(A\) stulpelių skaičius lygus \(B\) eilučių skaičiui.
\end{savybe}
\begin{teorema}
Jeigu sistemos
\[
\begin{cases}
ax+by=e,\\
cx+dy=f
\end{cases}
\]
koeficientų determinantas \(D=ad-bc\ne0\), sistema turi vienintelį sprendinį.
\end{teorema}
\begin{proof}
Geometriškai kiekviena lygtis aprašo tiesę. Kai \(D\ne0\), tiesės nėra lygiagrečios ir nesutampa, todėl jos kertasi viename taške.
\end{proof}

\subsection{Taisyklės ir algoritmai}
\begin{enumerate}[label=\arabic*.]
\item Sistemos koeficientus surašykite tvarkingai, laikydamiesi nežinomųjų eiliškumo.
\item Spręsdami Gauso metodu, naudokite elementariuosius eilučių veiksmus: sukeiskite eilutes, dauginkite eilutę iš nelygaus nuliui skaičiaus, prie eilutės pridėkite kitos eilutės kartotinį.
\item Siekite laiptuotos formos, tada atlikite atgalinį įstatymą.
\item Patikrinkite sprendinį pradinėje sistemoje.
\end{enumerate}

\subsection{Iliustruojantys pavyzdžiai}
\begin{enumerate}[label=\arabic*.]
\item Apskaičiuokime sandaugą:
\[
\begin{pmatrix}1&2\\3&4\end{pmatrix}
\begin{pmatrix}5\\6\end{pmatrix}
=
\begin{pmatrix}1\cdot5+2\cdot6\\3\cdot5+4\cdot6\end{pmatrix}
=
\begin{pmatrix}17\\39\end{pmatrix}.
\]
\textbf{Atsakymas:} \(\begin{pmatrix}17\\39\end{pmatrix}\).

\item Išspręskime sistemą
\[
\begin{cases}
2x+y=7,\\
x-y=2.
\end{cases}
\]
Sudėję lygtis gauname \(3x=9\), todėl \(x=3\). Tada \(3-y=2\), todėl \(y=1\).
\textbf{Atsakymas:} \(x=3,\ y=1\).

\item Ištirkime sistemą pagal determinantą:
\[
\begin{cases}
3x+2y=1,\\
6x+4y=5.
\end{cases}
\]
\[
D=\begin{vmatrix}3&2\\6&4\end{vmatrix}=12-12=0.
\]
Kairės pusės proporcingos, bet dešinės pusės ne: antroji lygtis nėra pirmosios kartotinis. Sistema sprendinių neturi.
\end{enumerate}

\subsection{Dažnos klaidos}
\begin{itemize}
\item Keičiant eilučių veiksmus pakeičiama tik viena lygties pusė.
\item Matricų sandaugoje dauginami atitinkami elementai kaip sudėtyje. Teisingai: eilutė dauginama iš stulpelio.
\item Determinantas \(ad-bc\) painiojamas su \(ab-cd\).
\end{itemize}

\subsection{Mini praktika}
\begin{enumerate}[label=\alph*)]
\item Apskaičiuokite \(\begin{vmatrix}4&-1\\2&3\end{vmatrix}\).
\item Išspręskite sistemą \(x+2y=8,\ 3x-y=1\).
\item Ar galima sudauginti \(2\times3\) ir \(3\times1\) matricas? Koks bus sandaugos matmuo?
\end{enumerate}

\subsection{Ryšys su kitomis temomis}
Matricos kompaktiškai užrašo sistemas, transformacijas ir duomenų lenteles. Jos padeda suprasti regresijos skaičiavimus, geometrijos transformacijas ir optimizavimo modelius.

\section{Optimizavimas ir modeliavimas}\label{sec:optimizavimas-ir-modeliavimas}

\begin{tcolorbox}[title=Primename iš ankstesnių klasių, colback=yellow!10]
Šiai temai reikia mokėti: \textbf{išvestinės taikymai}, nagrinėta 11 klasėje.
Jei reikia, pasikartok: \ref{sec:isvestines-taikymai}.
\end{tcolorbox}

\subsection{Apibrėžtys}
\begin{apibreztis}
Optimizavimo uždavinys yra uždavinys, kuriame reikia rasti didžiausią arba mažiausią dydžio reikšmę, laikantis duotų apribojimų.
\end{apibreztis}
\begin{apibreztis}
Matematinis modelis yra realios situacijos aprašymas kintamaisiais, lygtimis, nelygybėmis, funkcijomis arba duomenimis.
\end{apibreztis}
\begin{apibreztis}
Kritinis taškas yra apibrėžimo srities taškas, kuriame \(f'(x)=0\) arba išvestinė neegzistuoja.
\end{apibreztis}

\subsection{Teoremos ir savybės}
\begin{teorema}
Jeigu diferencijuojama funkcija turi ekstremumą vidiniame intervalo taške, tame taške jos išvestinė lygi nuliui.
\end{teorema}
\begin{proof}
Ekstremumo taške funkcijos reikšmės į abi puses pirmuoju artiniu nebedidėja ta pačia kryptimi. Todėl liestinės krypties koeficientas turi būti \(0\), t. y. \(f'(x)=0\).
\end{proof}
\begin{savybe}
Optimizuojant uždarame intervale būtina tikrinti ir kritinius taškus, ir intervalo galus.
\end{savybe}
\begin{savybe}
Jeigu \(f'\) keičia ženklą iš \(+\) į \(-\), turime maksimumą; jei iš \(-\) į \(+\), turime minimumą.
\end{savybe}

\subsection{Taisyklės ir algoritmai}
\begin{enumerate}[label=\arabic*.]
\item Apibrėžkite kintamuosius ir nurodykite jų leistinąsias reikšmes.
\item Sudarykite optimizuojamą funkciją vienu kintamuoju.
\item Raskite išvestinę ir kritinius taškus.
\item Palyginkite funkcijos reikšmes kritiniuose taškuose ir intervalo galuose.
\item Atsakymą parašykite kontekste: su vienetais ir patikrintu realumu.
\end{enumerate}

\subsection{Iliustruojantys pavyzdžiai}
\begin{enumerate}[label=\arabic*.]
\item Raskime \(f(x)=x^3-6x^2+9x+1\) ekstremumus.
\begin{align}
f'(x)&=3x^2-12x+9=3(x-1)(x-3).
\end{align}
Ženklų analizė: \(f'>0\), kai \(x<1\); \(f'<0\), kai \(1<x<3\); \(f'>0\), kai \(x>3\). Todėl ties \(x=1\) yra maksimumas, ties \(x=3\) -- minimumas.
\[
f(1)=5,\qquad f(3)=1.
\]
\textbf{Atsakymas:} maksimumas \(5\), minimumas \(1\).

\item Stačiakampio perimetras \(40\). Kokie matmenys duoda didžiausią plotą?
\[
2x+2y=40,\qquad y=20-x.
\]
\[
S(x)=x(20-x)=20x-x^2,\qquad 0<x<20.
\]
\[
S'(x)=20-2x=0,\qquad x=10,\qquad y=10.
\]
\textbf{Atsakymas:} didžiausią plotą duoda kvadratas \(10\times10\), plotas \(100\).

\item Iš \(24\) cm ilgio vielos lankstomas stačiakampis, kuris sukamas apie vieną kraštinę \(x\). Raskime didžiausią cilindro tūrį.
\[
2x+2r=24,\qquad r=12-x,\qquad V(x)=\pi r^2x=\pi(12-x)^2x.
\]
\[
V'(x)=\pi\left((12-x)^2-2x(12-x)\right)=\pi(12-x)(12-3x).
\]
Leistiname intervale \(0<x<12\) gauname \(x=4\), \(r=8\).
\textbf{Atsakymas:} didžiausias tūris \(V=256\pi\), kai aukštis \(4\) cm, spindulys \(8\) cm.
\end{enumerate}

\subsection{Dažnos klaidos}
\begin{itemize}
\item Sudaroma funkcija dviem kintamaisiais ir iš karto diferencijuojama, nepanaudojus apribojimo.
\item Rastas kritinis taškas laikomas atsakymu, bet nepatikrinami intervalo galai.
\item Atsakymas pateikiamas be vienetų arba neatitinka konteksto, pavyzdžiui, neigiamas ilgis.
\end{itemize}

\subsection{Mini praktika}
\begin{enumerate}[label=\alph*)]
\item Raskite \(f(x)=x^2-8x+10\) mažiausią reikšmę.
\item Skaičių \(30\) išskaidykite į du teigiamus dėmenis taip, kad jų sandauga būtų didžiausia.
\item Raskite funkcijos \(f(x)=x^3-3x\) didžiausią ir mažiausią reikšmes intervale \([-2;2]\).
\end{enumerate}

\subsection{Ryšys su kitomis temomis}
Optimizavimas sujungia funkcijų tyrimą, geometriją, integralų taikymus ir modeliavimą. Egzamine svarbiausia ne tik rasti išvestinę, bet ir teisingai sukurti funkciją iš teksto.

\section{Stereometrijos uždaviniai}\label{sec:stereometrijos-uzdaviniai}

\begin{tcolorbox}[title=Primename iš ankstesnių klasių, colback=yellow!10]
Šiai temai reikia mokėti: \textbf{erdvės geometrija}, nagrinėta 11 klasėje.
Jei reikia, pasikartok: \ref{sec:erdves-geometrija}.
\end{tcolorbox}

\subsection{Apibrėžtys}
\begin{apibreztis}
Stereometrija tiria erdvines figūras, jų tarpusavio padėtį, kampus, atstumus, pjūvius, plotus ir tūrius.
\end{apibreztis}
\begin{apibreztis}
Kampo tarp tiesės ir plokštumos dydis yra kampas tarp tiesės ir jos ortogonalios projekcijos į plokštumą.
\end{apibreztis}
\begin{apibreztis}
Dvisienis kampas yra kampas tarp dviejų pusplokštumių, turinčių bendrą kraštinę. Jis matuojamas kampu pjūvyje, statmename bendrai kraštinei.
\end{apibreztis}

\subsection{Teoremos ir savybės}
\begin{teorema}
Jeigu tiesė statmena dviem susikertančioms plokštumos tiesėms, ji statmena tai plokštumai.
\end{teorema}
\begin{proof}
Dvi susikertančios tiesės nusako dvi nepriklausomas plokštumos kryptis. Tiesė, statmena abiem šioms kryptims, statmena kiekvienai per susikirtimo tašką einančiai plokštumos tiesei, todėl statmena plokštumai.
\end{proof}
\begin{savybe}
Atstumas nuo taško iki plokštumos matuojamas statmens ilgiu, o ne bet kurios pasvirosios ilgiu.
\end{savybe}
\begin{savybe}
Prasilenkiančių tiesių kampas lygus kampui tarp lygiagrečiai perkeltų susikertančių tiesių.
\end{savybe}

\subsection{Taisyklės ir algoritmai}
\begin{enumerate}[label=\arabic*.]
\item Nusibraižykite aiškų brėžinį ir pažymėkite duotus ilgius bei kampus.
\item Ieškokite plokščiojo pjūvio, kuriame matomas ieškomas dydis.
\item Pažymėkite statmenis, projekcijas ir trikampius.
\item Taikykite Pitagoro teoremą, trigonometrines funkcijas, panašumą arba erdvinių kūnų formules.
\item Atsakymą patikrinkite pagal ribas: kampai turi būti realūs, ilgiai teigiami.
\end{enumerate}

\subsection{Iliustruojantys pavyzdžiai}
\begin{enumerate}[label=\arabic*.]
\item Stačiakampio gretasienio briaunos yra \(2\), \(3\) ir \(6\). Raskime įstrižainę.
\[
d=\sqrt{2^2+3^2+6^2}=\sqrt{49}=7.
\]
\textbf{Atsakymas:} \(7\).

\item Taisyklingos keturkampės piramidės pagrindo kraštinė \(6\), aukštinė \(4\). Raskime šoninės briaunos ilgį.
Nuo pagrindo centro iki viršūnės atstumas:
\[
R=\frac{6\sqrt2}{2}=3\sqrt2.
\]
Šoninė briauna:
\[
l=\sqrt{4^2+(3\sqrt2)^2}=\sqrt{16+18}=\sqrt{34}.
\]
\textbf{Atsakymas:} \(\sqrt{34}\).

\item Tiesė sudaro su plokštuma \(30^\circ\) kampą. Pasviroji nuo taško iki plokštumos yra \(10\). Raskime atstumą nuo taško iki plokštumos.
\[
h=10\sin30^\circ=10\cdot\frac12=5.
\]
\textbf{Atsakymas:} \(5\).
\end{enumerate}

\subsection{Dažnos klaidos}
\begin{itemize}
\item Kampas tarp tiesės ir plokštumos painiojamas su kampu tarp tiesės ir bet kurios plokštumos tiesės.
\item Vietoj statmens iki plokštumos naudojama pasviroji.
\item Erdvinis uždavinys sprendžiamas be pjūvio, todėl į vieną trikampį sudedami taškai, kurie nėra vienoje plokštumoje.
\end{itemize}

\subsection{Mini praktika}
\begin{enumerate}[label=\alph*)]
\item Kubo briauna \(5\). Raskite jo erdvinę įstrižainę.
\item Taisyklingos trikampės piramidės aukštinė \(6\), o atstumas nuo pagrindo centro iki viršūnės \(2\sqrt3\). Raskite šoninę briauną.
\item Paaiškinkite, kodėl atstumas nuo taško iki plokštumos visada ne didesnis už bet kurios pasvirosios ilgį.
\end{enumerate}

\subsection{Ryšys su kitomis temomis}
Stereometrija remiasi planimetrija, vektoriais ir trigonometrija. Pjūviai dažnai paverčia erdvinį uždavinį į plokštumos geometrijos uždavinį.

\section{Tikimybiniai skirstiniai}\label{sec:tikimybiniai-skirstiniai}

\begin{tcolorbox}[title=Primename iš ankstesnių klasių, colback=yellow!10]
Šiai temai reikia mokėti: \textbf{sąlyginė tikimybė ir nepriklausomumas}, nagrinėta 11 klasėje.
Jei reikia, pasikartok: \ref{sec:salygine-tikimybe-ir-nepriklausomumas}.
\end{tcolorbox}

\subsection{Apibrėžtys}
\begin{apibreztis}
Atsitiktinio dydžio skirstinys nusako, kokias reikšmes dydis gali įgyti ir su kokiomis tikimybėmis.
\end{apibreztis}
\begin{apibreztis}
Matematinė viltis yra ilgalaikis atsitiktinio dydžio vidurkis:
\[
E(X)=\sum x_i p_i.
\]
\end{apibreztis}
\begin{apibreztis}
Dispersija matuoja sklaidos dydį:
\[
D(X)=\sum (x_i-E(X))^2p_i,\qquad \sigma(X)=\sqrt{D(X)}.
\]
\end{apibreztis}

\subsection{Teoremos ir savybės}
\begin{savybe}
Diskrečiojo skirstinio tikimybių suma lygi \(1\).
\end{savybe}
\begin{teorema}
Binominio skirstinio \(X\sim B(n,p)\) matematinė viltis yra \(E(X)=np\), o dispersija \(D(X)=np(1-p)\).
\end{teorema}
\begin{proof}
Binominis dydis yra \(n\) nepriklausomų Bernulio bandymų suma. Kiekvieno bandymo viltis lygi \(p\), dispersija \(p(1-p)\), todėl sumos viltis ir dispersija yra atitinkamų dydžių sumos.
\end{proof}
\begin{savybe}
Jeigu \(X\sim B(n,p)\), tai
\[
P(X=k)=\binom{n}{k}p^k(1-p)^{n-k}.
\]
\end{savybe}

\subsection{Taisyklės ir algoritmai}
\begin{enumerate}[label=\arabic*.]
\item Nustatykite, ką reiškia atsitiktinis dydis \(X\).
\item Išvardykite galimas reikšmes ir priskirkite tikimybes.
\item Patikrinkite, ar tikimybių suma lygi \(1\).
\item Skaičiuokite viltį, dispersiją, standartinį nuokrypį arba konkrečias tikimybes.
\item Jei yra \(n\) vienodų nepriklausomų bandymų su dviem baigtimis, patikrinkite, ar tinka binominis modelis.
\end{enumerate}

\subsection{Iliustruojantys pavyzdžiai}
\begin{enumerate}[label=\arabic*.]
\item Dydis \(X\) įgyja reikšmes \(0,1,2\) su tikimybėmis \(0{,}2,0{,}5,0{,}3\). Raskime \(E(X)\).
\[
E(X)=0\cdot0{,}2+1\cdot0{,}5+2\cdot0{,}3=1{,}1.
\]
\textbf{Atsakymas:} \(1{,}1\).

\item Raskime to paties dydžio dispersiją.
\begin{align}
D(X)&=(0-1{,}1)^2\cdot0{,}2+(1-1{,}1)^2\cdot0{,}5+(2-1{,}1)^2\cdot0{,}3\\
&=1{,}21\cdot0{,}2+0{,}01\cdot0{,}5+0{,}81\cdot0{,}3\\
&=0{,}49.
\end{align}
\textbf{Atsakymas:} \(D(X)=0{,}49\), \(\sigma=0{,}7\).

\item Metama moneta \(5\) kartus. Raskime tikimybę gauti lygiai \(3\) herbus.
\[
P(X=3)=\binom53\left(\frac12\right)^3\left(\frac12\right)^2=\frac{10}{32}=\frac{5}{16}.
\]
\textbf{Atsakymas:} \(\frac{5}{16}\).
\end{enumerate}

\subsection{Dažnos klaidos}
\begin{itemize}
\item Viltis painiojama su labiausiai tikėtina reikšme. Ji gali net nebūti galima reikšmė.
\item Skirstinio lentelėje tikimybių suma nėra \(1\).
\item Binominis modelis taikomas priklausomiems bandymams be pagrindo.
\end{itemize}

\subsection{Mini praktika}
\begin{enumerate}[label=\alph*)]
\item \(X\) įgyja \(1,2,3\) su tikimybėmis \(0{,}4,0{,}4,0{,}2\). Raskite \(E(X)\).
\item Jei \(X\sim B(12,0{,}25)\), raskite \(E(X)\) ir \(D(X)\).
\item Raskite \(P(X=2)\), kai \(X\sim B(4,0{,}5)\).
\end{enumerate}

\subsection{Ryšys su kitomis temomis}
Skirstiniai remiasi kombinatorika ir tikimybių savybėmis. Normalusis skirstinys yra tolydus atsitiktinių dydžių modelis, todėl siejasi su funkcijų grafikais ir plotais po kreive.

\section{Statistinė išvada ir hipotezės}\label{sec:statistine-isvada-ir-hipotezes}

\begin{tcolorbox}[title=Primename iš ankstesnių klasių, colback=yellow!10]
Šiai temai reikia mokėti: \textbf{aprašomoji statistika ir normalusis modelis}, nagrinėta 11 klasėje.
Jei reikia, pasikartok: \ref{sec:aprasomoji-statistika-ir-normalusis-modelis}.
\end{tcolorbox}

\subsection{Apibrėžtys}
\begin{apibreztis}
Populiacija yra visa tiriama objektų ar žmonių aibė, o imtis -- jos dalis, pagal kurią daroma išvada apie populiaciją.
\end{apibreztis}
\begin{apibreztis}
Nulinė hipotezė \(H_0\) yra tikrinamas pradinis teiginys, o alternatyvioji hipotezė \(H_1\) nurodo, kokio pokyčio ar skirtumo ieškoma.
\end{apibreztis}
\begin{apibreztis}
\(p\)-reikšmė yra tikimybė gauti bent tokį kraštutinį rezultatą, jeigu \(H_0\) yra teisinga.
\end{apibreztis}

\subsection{Teoremos ir savybės}
\begin{savybe}
Didesnė atsitiktinė imtis paprastai mažina įverčio paklaidą, nes atsitiktiniai svyravimai labiau išsilygina.
\end{savybe}
\begin{teorema}
Jeigu imtis didelė ir stebėjimai nepriklausomi, imties vidurkio skirstinys dažnai artimas normaliajam.
\end{teorema}
\begin{proof}
Tai centrinės ribinės teoremos idėja: daug nepriklausomų mažų įtakų suma arba vidurkis turi stabilų, varpo formos skirstinį.
\end{proof}
\begin{savybe}
Statistinė išvada visada priklauso nuo duomenų rinkimo būdo. Neatsitiktinė arba šališka imtis gali duoti klaidingą išvadą net tada, kai skaičiavimai atlikti teisingai.
\end{savybe}

\subsection{Taisyklės ir algoritmai}
\begin{enumerate}[label=\arabic*.]
\item Aiškiai apibrėžkite populiaciją, imtį ir matuojamą kintamąjį.
\item Suformuluokite \(H_0\) ir \(H_1\).
\item Pasirinkite reikšmingumo lygmenį \(\alpha\), dažnai \(\alpha=0{,}05\).
\item Apskaičiuokite testinę statistiką arba \(p\)-reikšmę.
\item Jeigu \(p\le\alpha\), atmeskite \(H_0\); jeigu \(p>\alpha\), nepakanka pagrindo atmesti \(H_0\).
\item Sprendimą išsakykite kontekste, vengdami per stiprių teiginių.
\end{enumerate}

\subsection{Iliustruojantys pavyzdžiai}
\begin{enumerate}[label=\arabic*.]
\item Imties vidurkis \(52\), hipotezinis vidurkis \(50\), standartinė paklaida \(1\). Raskime \(z\).
\[
z=\frac{52-50}{1}=2.
\]
\textbf{Atsakymas:} \(z=2\).

\item Apklausus \(100\) mokinių, \(60\) pasirinko A. Raskime proporcijos įvertį ir apytikslę standartinę paklaidą.
\[
\hat p=\frac{60}{100}=0{,}6,\qquad
SE\approx\sqrt{\frac{0{,}6\cdot0{,}4}{100}}=\sqrt{0{,}0024}\approx0{,}049.
\]
\textbf{Atsakymas:} \(\hat p=0{,}6\), standartinė paklaida apie \(0{,}049\).

\item Jei pasikliautinojo intervalo centras \(10\), paklaida \(0{,}8\), užrašykime intervalą.
\[
I=[10-0{,}8;\,10+0{,}8]=[9{,}2;\,10{,}8].
\]
\textbf{Atsakymas:} \([9{,}2;\,10{,}8]\).
\end{enumerate}

\subsection{Dažnos klaidos}
\begin{itemize}
\item Teigiama, kad \(H_0\) ``įrodyta'', kai jos neatmetame. Teisingiau: duomenys nesuteikė pakankamai pagrindo ją atmesti.
\item Painiojama koreliacija ir priežastingumas.
\item Iš šališkos apklausos daroma išvada apie visą populiaciją.
\end{itemize}

\subsection{Mini praktika}
\begin{enumerate}[label=\alph*)]
\item Suformuluokite \(H_0\) ir \(H_1\), jei tiriama, ar vidutinis kelionės į mokyklą laikas viršija \(25\) min.
\item Paaiškinkite, ką reiškia \(p=0{,}03\), kai \(\alpha=0{,}05\).
\item Kodėl savanoriška internetinė apklausa dažnai nėra gera imtis visai mokyklai?
\end{enumerate}

\subsection{Ryšys su kitomis temomis}
Statistinė išvada remiasi duomenų aprašymu, tikimybiniais skirstiniais ir modeliavimu. Ji padeda pereiti nuo skaičiavimo prie pagrįsto sprendimo, bet reikalauja kritiškai vertinti duomenų kilmę.

\section{Egzamino sintezė ir problemų sprendimas}\label{sec:egzamino-sinteze-ir-problemu-sprendimas}

\subsection{Apibrėžtys}
\begin{apibreztis}
Sintezės uždavinys yra uždavinys, kuriame reikia sujungti kelias matematikos sritis ir pasirinkti tinkamą sprendimo strategiją.
\end{apibreztis}
\begin{apibreztis}
Sprendimo strategija yra sąmoningas veiksmų planas: suprasti sąlygą, pasirinkti modelį, atlikti skaičiavimus ir patikrinti rezultatą.
\end{apibreztis}
\begin{apibreztis}
Pilnas sprendimas yra ne tik atsakymas, bet ir argumentuotas kelias iki jo: apibrėžti kintamieji, formulės, pertvarkymai, skaičiavimai ir išvada.
\end{apibreztis}

\subsection{Teoremos ir savybės}
\begin{savybe}
Sudėtingame uždavinyje tarpiniai rezultatai turi būti tikrinami pagal prasmę, vienetus, ženklą ir leistinąsias reikšmes.
\end{savybe}
\begin{savybe}
Tas pats uždavinys gali turėti kelis sprendimo kelius: algebrinį, grafinį, geometrinį, tikimybinį arba skaitinį.
\end{savybe}
\begin{savybe}
Brandos egzamino užduotyse ypač vertinamas ne vien galutinis skaičius, bet ir argumentavimas: kodėl pasirinktas metodas tinka ir kodėl gautas rezultatas atsako į klausimą.
\end{savybe}

\subsection{Taisyklės ir algoritmai}
\begin{enumerate}[label=\arabic*.]
\item Perskaitykite sąlygą du kartus: pirmą kartą dėl situacijos, antrą kartą dėl duomenų.
\item Išskirkite duotus dydžius, ieškomą dydį, apribojimus ir vienetus.
\item Atpažinkite temas: funkcijos, trigonometrija, geometrija, tikimybės, statistika, integralai.
\item Sudarykite modelį ir spręskite nuosekliai.
\item Patikrinkite atsakymą kitu būdu arba bent įvertinkite jo realumą.
\item Galutinį atsakymą parašykite aiškiu sakiniu.
\end{enumerate}

\subsection{Iliustruojantys pavyzdžiai}
\begin{enumerate}[label=\arabic*.]
\item Funkcija \(f(x)=x^2-6x+8\). Raskime jos nulius ir mažiausią reikšmę.
\[
x^2-6x+8=(x-2)(x-4),
\]
todėl nuliai \(x=2\) ir \(x=4\). Mažiausia reikšmė yra viršūnėje:
\[
x=\frac{2+4}{2}=3,\qquad f(3)=9-18+8=-1.
\]
\textbf{Atsakymas:} nuliai \(2\) ir \(4\), mažiausia reikšmė \(-1\).

\item Dėžėje yra \(3\) raudoni ir \(2\) mėlyni rutuliai. Traukiami \(2\) be grąžinimo. Raskime tikimybę ištraukti bent vieną mėlyną.
Patogiau skaičiuoti priešingą įvykį: abu raudoni.
\[
P(\text{abu raudoni})=\frac35\cdot\frac24=\frac3{10}.
\]
\[
P(\text{bent vienas mėlynas})=1-\frac3{10}=\frac7{10}.
\]
\textbf{Atsakymas:} \(\frac7{10}\).

\item Stačiojo trikampio statiniai \(6\) ir \(8\). Jo plotas sukamas apie statinį \(6\). Raskime susidariusio kūgio tūrį.
Sukant apie statinį \(6\), kūgio aukštis \(h=6\), spindulys \(r=8\).
\[
V=\frac13\pi r^2h=\frac13\pi\cdot64\cdot6=128\pi.
\]
\textbf{Atsakymas:} \(128\pi\).
\end{enumerate}

\subsection{Dažnos klaidos}
\begin{itemize}
\item Atsakoma į ne tą klausimą: pavyzdžiui, randamas spindulys, nors prašoma tūrio.
\item Per anksti apvalinama, todėl vėlesni skaičiavimai praranda tikslumą.
\item Sprendime nėra paaiškinta, kodėl pasirinkta formulė tinka.
\end{itemize}

\subsection{Mini praktika}
\begin{enumerate}[label=\alph*)]
\item Sudarykite trijų žingsnių planą uždaviniui: rasti plotą tarp dviejų funkcijų grafikų.
\item Kodėl tikimybių uždavinyje naudinga ieškoti priešingo įvykio?
\item Funkcija \(g(x)=x^3-3x\). Kokias temas reikėtų prisiminti, norint nubraižyti jos grafiko eskizą?
\end{enumerate}

\subsection{Ryšys su kitomis temomis}
Egzamino sintezė apibendrina visą kursą: skaičius, algebrą, funkcijas, geometriją, statistiką ir tikimybę. Gera strategija dažnai svarbesnė už ilgą skaičiavimą.

\section{Trigonometrinės lygtys ir nelygybės}\label{sec:trigonometrines-lygtys-ir-nelygybes}

\subsection{Apibrėžtys}
\begin{apibreztis}
Trigonometrinė lygtis yra lygtis, kurioje nežinomasis yra trigonometrinės funkcijos argumente, pavyzdžiui, \(\sin x=\frac12\), \(\cos 2x=0\), \(\tan x>1\).
\end{apibreztis}
\begin{apibreztis}
Bendrasis trigonometrinės lygties sprendinys aprašo visas reikšmes, besikartojančias dėl trigonometrinių funkcijų periodiškumo.
\end{apibreztis}
\begin{apibreztis}
Trigonometrinė nelygybė sprendžiama nustatant, kuriuose vienetinio apskritimo arba grafiko intervaluose funkcijos reikšmės tenkina nelygybę.
\end{apibreztis}

\subsection{Teoremos ir savybės}
\begin{savybe}
Pagrindiniai periodai:
\[
\sin(x+2\pi)=\sin x,\qquad \cos(x+2\pi)=\cos x,\qquad \tan(x+\pi)=\tan x.
\]
\end{savybe}
\begin{savybe}
Pagrindinės lygčių formulės:
\[
\sin x=a\Rightarrow x=(-1)^n\arcsin a+n\pi,\quad |a|\le1,
\]
\[
\cos x=a\Rightarrow x=\pm\arccos a+2\pi n,\quad |a|\le1,
\]
\[
\tan x=a\Rightarrow x=\arctan a+\pi n,\quad n\in\mathbb{Z}.
\]
\end{savybe}
\begin{savybe}
Dažnai naudojamos tapatybės:
\[
\sin^2x+\cos^2x=1,\qquad
\sin2x=2\sin x\cos x,\qquad
\cos2x=\cos^2x-\sin^2x.
\]
Išplėstiniame kurse taip pat taikomos sumos ir skirtumo formulės.
\end{savybe}

\subsection{Taisyklės ir algoritmai}
\begin{enumerate}[label=\arabic*.]
\item Pertvarkykite lygtį iki vienos pagrindinės funkcijos, jei įmanoma.
\item Jei yra sandauga, naudokite nulio sandaugos taisyklę.
\item Užrašykite bendruosius sprendinius su \(n\in\mathbb{Z}\).
\item Jei prašoma sprendinių intervale, iš bendrųjų formulių atrinkite tik į intervalą patenkančius sprendinius.
\item Nelygybei naudokite vienetinį apskritimą arba grafiką ir tik tada pridėkite periodus.
\end{enumerate}

\subsection{Iliustruojantys pavyzdžiai}
\begin{enumerate}[label=\arabic*.]
\item Išspręskime \(\sin x=\frac12\).
Viename periode sprendiniai yra \(\frac{\pi}{6}\) ir \(\frac{5\pi}{6}\). Todėl
\[
x=\frac{\pi}{6}+2\pi n\quad\text{arba}\quad x=\frac{5\pi}{6}+2\pi n,\qquad n\in\mathbb{Z}.
\]
\textbf{Atsakymas:} \(x=\frac{\pi}{6}+2\pi n\) arba \(x=\frac{5\pi}{6}+2\pi n\).

\item Raskime \(\cos 2x=0\) sprendinius intervale \([0;2\pi]\).
\[
2x=\frac{\pi}{2}+\pi n,\qquad x=\frac{\pi}{4}+\frac{\pi n}{2}.
\]
Atrenkame:
\[
x=\frac{\pi}{4},\ \frac{3\pi}{4},\ \frac{5\pi}{4},\ \frac{7\pi}{4}.
\]
\textbf{Atsakymas:} \(\frac{\pi}{4},\frac{3\pi}{4},\frac{5\pi}{4},\frac{7\pi}{4}\).

\item Išspręskime \(\sin x\ge\frac{\sqrt3}{2}\).
Viename periode \([0;2\pi]\) sinusas toks didelis nuo \(\frac{\pi}{3}\) iki \(\frac{2\pi}{3}\). Todėl
\[
x\in\left[\frac{\pi}{3}+2\pi n;\frac{2\pi}{3}+2\pi n\right],\qquad n\in\mathbb{Z}.
\]
\textbf{Atsakymas:} nurodyti intervalai.
\end{enumerate}

\subsection{Dažnos klaidos}
\begin{itemize}
\item Pamirštamas periodas ir pateikiamas tik vienas sprendinys.
\item Sprendžiant \(\cos x=a\), prarandamas antras sprendinys \(-\arccos a+2\pi n\).
\item Nelygybė sprendžiama kaip lygtis, nepateikiant intervalų.
\item Dalijama iš \(\sin x\) arba \(\cos x\), nepatikrinus atvejų, kai šis daugiklis lygus nuliui.
\end{itemize}

\subsection{Mini praktika}
\begin{enumerate}[label=\alph*)]
\item Išspręskite \(\cos x=-\frac12\).
\item Raskite \(\tan x=1\) sprendinius intervale \([0;2\pi]\).
\item Išspręskite \(\cos x<0\).
\end{enumerate}

\subsection{Ryšys su kitomis temomis}
Trigonometrinės lygtys remiasi vienetiniu apskritimu, funkcijų grafikais ir tapatybėmis. Jos vėliau naudojamos išvestinių, integralų, periodinių procesų ir geometrijos uždaviniuose.

\section{Išvestinė, tolydumas ir liestinė}\label{sec:isvestine-tolydumas-ir-liestine}

\subsection{Apibrėžtys}
\begin{apibreztis}
Funkcija yra tolydi taške \(x=a\), jeigu jos reikšmės artėjant prie \(a\) artėja prie \(f(a)\). Intuityviai grafikas ties tuo tašku neturi šuolio ar skylės.
\end{apibreztis}
\begin{apibreztis}
Funkcijos išvestinė taške \(a\) yra riba
\[
f'(a)=\lim_{h\to0}\frac{f(a+h)-f(a)}{h},
\]
jeigu ši riba egzistuoja.
\end{apibreztis}
\begin{apibreztis}
Liestinė funkcijos grafikui taške \(x=a\) yra tiesė, kurios krypties koeficientas lygus \(f'(a)\). Jos lygtis:
\[
y=f(a)+f'(a)(x-a).
\]
\end{apibreztis}

\subsection{Teoremos ir savybės}
\begin{savybe}
Jeigu funkcija diferencijuojama taške, tai ji tame taške yra tolydi. Atvirkščias teiginys nebūtinai teisingas.
\end{savybe}
\begin{savybe}
Išvestinių taisyklės:
\[
(cf)'=cf',\quad (f+g)'=f'+g',\quad (fg)'=f'g+fg',
\]
\[
\left(\frac{f}{g}\right)'=\frac{f'g-fg'}{g^2},\quad g\ne0,\qquad
(f(g(x)))'=f'(g(x))g'(x).
\]
\end{savybe}
\begin{savybe}
Elementariųjų funkcijų išvestinės:
\[
(x^n)'=nx^{n-1},\quad (e^x)'=e^x,\quad (\ln x)'=\frac1x,
\]
\[
(\sin x)'=\cos x,\quad (\cos x)'=-\sin x,\quad (\tan x)'=\frac{1}{\cos^2x}.
\]
\end{savybe}

\subsection{Taisyklės ir algoritmai}
\begin{enumerate}[label=\arabic*.]
\item Nustatykite funkcijos apibrėžimo sritį ir galimus netolydumo taškus.
\item Išvestinei rinkitės tinkamą taisyklę: suma, sandauga, dalmuo, sudėtinė funkcija.
\item Liestinei raskite \(f(a)\) ir \(f'(a)\).
\item Funkcijos tyrimui išspręskite \(f'(x)=0\) ir sudarykite išvestinės ženklų lentelę.
\item Grafiko eskizui pažymėkite nulius, ekstremumus, didėjimo ir mažėjimo intervalus.
\end{enumerate}

\subsection{Iliustruojantys pavyzdžiai}
\begin{enumerate}[label=\arabic*.]
\item Raskime \(f(x)=x^3-4x\) išvestinę ir kritinius taškus.
\[
f'(x)=3x^2-4.
\]
\[
3x^2-4=0,\qquad x=\pm\frac{2}{\sqrt3}=\pm\frac{2\sqrt3}{3}.
\]
\textbf{Atsakymas:} kritiniai taškai \(x=\pm\frac{2\sqrt3}{3}\).

\item Raskime funkcijos \(f(x)=x^2+1\) liestinės lygtį taške \(x=2\).
\[
f(2)=5,\qquad f'(x)=2x,\qquad f'(2)=4.
\]
\[
y=5+4(x-2)=4x-3.
\]
\textbf{Atsakymas:} \(y=4x-3\).

\item Raskime \(g(x)=\sin(3x^2)\) išvestinę.
\[
g'(x)=\cos(3x^2)\cdot6x=6x\cos(3x^2).
\]
\textbf{Atsakymas:} \(6x\cos(3x^2)\).
\end{enumerate}

\subsection{Dažnos klaidos}
\begin{itemize}
\item Sudėtinės funkcijos išvestinėje pamirštamas vidinės funkcijos daugiklis.
\item Liestinės lygtyje vietoj taško \((a;f(a))\) naudojamas \((a;f'(a))\).
\item Kritinis taškas sutapatinamas su ekstremumu, nors būtina patikrinti išvestinės ženklo kitimą.
\end{itemize}

\subsection{Mini praktika}
\begin{enumerate}[label=\alph*)]
\item Raskite \(f(x)=\frac{x^2+1}{x}\) išvestinę.
\item Parašykite \(y=\ln x\) liestinės lygtį taške \(x=1\).
\item Ištirkite, kur \(f(x)=x^3-12x\) didėja ir mažėja.
\end{enumerate}

\subsection{Ryšys su kitomis temomis}
Išvestinė yra pagrindinis funkcijų tyrimo ir optimizavimo įrankis. Ji taip pat paruošia integralą, nes pirmykštė funkcija apibrėžiama per išvestinę.

\section{Briaunainiai, sukiniai ir pjūviai}\label{sec:briaunainiai-sukiniai-ir-pjuviai}

\subsection{Apibrėžtys}
\begin{apibreztis}
Briaunainis yra erdvinis kūnas, apribotas daugiakampiais. Prizmės ir piramidės yra svarbiausi mokykliniai briaunainiai.
\end{apibreztis}
\begin{apibreztis}
Sukinys yra kūnas, gaunamas sukant plokštumos figūrą apie tiesę. Ritinys, kūgis, nupjautinis kūgis, sfera ir rutulys yra sukiniai.
\end{apibreztis}
\begin{apibreztis}
Pjūvis yra figūra, gaunama erdvinį kūną kertant plokštuma. Ašinis pjūvis eina per sukinio ašį.
\end{apibreztis}

\subsection{Teoremos ir savybės}
\begin{savybe}
Stačiosios prizmės tūris \(V=S_p h\), kur \(S_p\) -- pagrindo plotas, \(h\) -- aukštis.
\end{savybe}
\begin{savybe}
Piramidės tūris \(V=\frac13 S_p h\), kūgio tūris \(V=\frac13\pi r^2h\), ritinio tūris \(V=\pi r^2h\).
\end{savybe}
\begin{savybe}
Rutulio tūris \(V=\frac43\pi r^3\), paviršiaus plotas \(S=4\pi r^2\).
\end{savybe}
\begin{savybe}
Pjūvis, lygiagretus pagrindui, dažnai yra panašus į pagrindą. Ilgiai keičiasi pagal panašumo koeficientą, plotai -- pagal jo kvadratą.
\end{savybe}

\subsection{Taisyklės ir algoritmai}
\begin{enumerate}[label=\arabic*.]
\item Atpažinkite kūną ir jo elementus: pagrindą, aukštinę, apotemą, sudaromąją, spindulį.
\item Jei reikia pjūvio, nusibraižykite būtent tą plokštumos figūrą.
\item Pasirinkite formulę tūriui, šoniniam paviršiui arba visam paviršiui.
\item Trūkstamus ilgius raskite iš stačiųjų trikampių, panašumo arba trigonometrinių funkcijų.
\item Atskirai nurodykite, ar skaičiuojamas šoninis, ar visas paviršius.
\end{enumerate}

\subsection{Iliustruojantys pavyzdžiai}
\begin{enumerate}[label=\arabic*.]
\item Ritinio spindulys \(3\), aukštis \(7\). Raskime tūrį ir šoninį paviršių.
\[
V=\pi r^2h=\pi\cdot9\cdot7=63\pi.
\]
\[
S_{\text{šon}}=2\pi rh=2\pi\cdot3\cdot7=42\pi.
\]
\textbf{Atsakymas:} \(V=63\pi\), \(S_{\text{šon}}=42\pi\).

\item Kūgio spindulys \(5\), aukštis \(12\). Raskime sudaromąją ir šoninį paviršių.
\[
l=\sqrt{5^2+12^2}=13,\qquad S_{\text{šon}}=\pi rl=65\pi.
\]
\textbf{Atsakymas:} \(l=13\), \(S_{\text{šon}}=65\pi\).

\item Taisyklingos keturkampės piramidės pagrindo kraštinė \(8\), aukštis \(3\). Raskime tūrį.
\[
S_p=8^2=64,\qquad V=\frac13\cdot64\cdot3=64.
\]
\textbf{Atsakymas:} \(64\).
\end{enumerate}

\subsection{Dažnos klaidos}
\begin{itemize}
\item Painiojama kūgio aukštinė ir sudaromoji.
\item Paviršiaus plote pamirštamas vienas arba abu pagrindai.
\item Pjūvio plotas keičiamas pagal ilgių koeficientą, nors plotams reikia koeficiento kvadrato.
\end{itemize}

\subsection{Mini praktika}
\begin{enumerate}[label=\alph*)]
\item Raskite rutulio, kurio spindulys \(2\), tūrį ir paviršiaus plotą.
\item Raskite kūgio tūrį, kai \(r=4\), \(h=9\).
\item Kuo skiriasi ritinio ašinis pjūvis nuo pjūvio, lygiagretaus pagrindui?
\end{enumerate}

\subsection{Ryšys su kitomis temomis}
Briaunainių ir sukinių uždaviniai remiasi stereometrija, trigonometrija ir funkcijomis. Sukinių tūriai išplėstiniame kurse natūraliai susiję su integralais.

\section{Regresija ir koreliacija}\label{sec:regresija-ir-koreliacija}

\subsection{Apibrėžtys}
\begin{apibreztis}
Koreliacija apibūdina dviejų kiekybinių kintamųjų tiesinio ryšio kryptį ir stiprumą.
\end{apibreztis}
\begin{apibreztis}
Koreliacijos koeficientas \(r\) tenkina \(-1\le r\le1\). Kai \(r>0\), ryšys teigiamas; kai \(r<0\), ryšys neigiamas; kai \(r\approx0\), tiesinis ryšys silpnas.
\end{apibreztis}
\begin{apibreztis}
Tiesinė regresija aprašo priklausomąjį kintamąjį \(y\) aiškinamuoju kintamuoju \(x\) tiesės modeliu \(y=ax+b\).
\end{apibreztis}
\begin{apibreztis}
Determinacijos koeficientas \(R^2\) parodo, kokią priklausomojo kintamojo sklaidos dalį paaiškina modelis.
\end{apibreztis}

\subsection{Teoremos ir savybės}
\begin{savybe}
Koreliacija neįrodo priežastingumo. Du kintamieji gali judėti kartu dėl trečio veiksnio arba atsitiktinumo.
\end{savybe}
\begin{savybe}
Regresijos tiesė tinka prognozuoti tik tame \(x\) reikšmių intervale, kuriame buvo rinkti duomenys. Prognozė toli už intervalo vadinama ekstrapoliacija ir yra rizikinga.
\end{savybe}
\begin{savybe}
Jeigu modelyje yra vienas aiškinamasis kintamasis, dažnai \(R^2=r^2\).
\end{savybe}

\subsection{Taisyklės ir algoritmai}
\begin{enumerate}[label=\arabic*.]
\item Nubraižykite taškų diagramą ir įvertinkite, ar ryšys panašus į tiesinį.
\item Nuspręskite, kuris kintamasis yra aiškinamasis \(x\), o kuris priklausomasis \(y\).
\item Skaitmeninėmis priemonėmis raskite regresijos tiesę \(y=ax+b\), koreliacijos koeficientą \(r\) ir \(R^2\).
\item Interpretuokite \(a\): kiek vidutiniškai keičiasi \(y\), kai \(x\) padidėja vienetu.
\item Patikrinkite, ar išvada neperžengia duomenų ribų ir nevadina koreliacijos priežastimi.
\end{enumerate}

\subsection{Iliustruojantys pavyzdžiai}
\begin{enumerate}[label=\arabic*.]
\item Duomenims gauta regresijos tiesė \(\hat y=1{,}8x+12\). Ką reiškia koeficientas \(1{,}8\)?
\textbf{Sprendimas.} Kai \(x\) padidėja vienu vienetu, modelis prognozuoja, kad \(y\) vidutiniškai padidėja \(1{,}8\) vieneto.

\item Jei \(r=-0{,}86\), apibūdinkime ryšį.
\textbf{Sprendimas.} Ryšys stiprus neigiamas: didėjant vienam kintamajam, kitas paprastai mažėja, bet tai savaime neįrodo priežasties.

\item Modelio \(R^2=0{,}72\). Interpretuokime.
\textbf{Sprendimas.} Apie \(72\%\) priklausomojo kintamojo sklaidos paaiškinama tiesiniu modeliu su pasirinktu aiškinamuoju kintamuoju.
\end{enumerate}

\subsection{Dažnos klaidos}
\begin{itemize}
\item Iš koreliacijos daroma priežastinė išvada be papildomo tyrimo.
\item Regresijos modelis taikomas už duomenų intervalo ribų.
\item Ignoruojami išskirtys, kurios gali smarkiai pakeisti tiesę ir koreliaciją.
\end{itemize}

\subsection{Mini praktika}
\begin{enumerate}[label=\alph*)]
\item Paaiškinkite, ką reiškia \(r=0{,}15\).
\item Regresijos tiesė \(\hat y=-2x+50\). Kokia prognozė, kai \(x=12\)?
\item Kodėl prieš skaičiuojant regresiją verta nubraižyti taškų diagramą?
\end{enumerate}

\subsection{Ryšys su kitomis temomis}
Regresija jungia funkcijas ir statistiką: tiesė tampa duomenų modeliu. Ji taip pat remiasi kritiniu mąstymu apie imtį, kintamuosius ir paklaidas.

\section{Klasikiniai ir neklasikiniai tikimybiniai modeliai}\label{sec:klasikiniai-ir-neklasikiniai-tikimybiniai-modeliai}

\subsection{Apibrėžtys}
\begin{apibreztis}
Klasikinis tikimybinis modelis taikomas tada, kai visos elementariosios baigtys vienodai tikėtinos. Tada
\[
P(A)=\frac{\text{palankių baigčių skaičius}}{\text{visų baigčių skaičius}}.
\]
\end{apibreztis}
\begin{apibreztis}
Neklasikiniame modelyje baigtys nebūtinai vienodai tikėtinos, todėl tikimybės nustatomos iš sąlygos, duomenų, medžio arba modelio.
\end{apibreztis}
\begin{apibreztis}
Įvykiai yra nesutaikomi, jei negali įvykti kartu. Įvykiai yra nepriklausomi, jei vieno įvykimas nekeičia kito tikimybės.
\end{apibreztis}
\begin{apibreztis}
Bernulio bandymas turi dvi baigtis: sėkmę su tikimybe \(p\) ir nesėkmę su tikimybe \(1-p\).
\end{apibreztis}

\subsection{Teoremos ir savybės}
\begin{savybe}
\[
P(\Omega)=1,\qquad P(\varnothing)=0,\qquad P(\overline A)=1-P(A).
\]
\end{savybe}
\begin{savybe}
Jeigu \(A\) ir \(B\) nesutaikomi, tai
\[
P(A\cup B)=P(A)+P(B).
\]
Bendruoju atveju
\[
P(A\cup B)=P(A)+P(B)-P(A\cap B).
\]
\end{savybe}
\begin{savybe}
Jeigu \(A\) ir \(B\) nepriklausomi, tai
\[
P(A\cap B)=P(A)P(B).
\]
\end{savybe}
\begin{savybe}
Bernulio bandymuose tikimybė gauti tiksliai \(k\) sėkmių per \(n\) bandymų:
\[
P(X=k)=\binom nk p^k(1-p)^{n-k}.
\]
\end{savybe}

\subsection{Taisyklės ir algoritmai}
\begin{enumerate}[label=\arabic*.]
\item Nuspręskite, ar baigtys vienodai tikėtinos. Jei taip, galima taikyti klasikinę formulę.
\item Sudėtiniam bandymui braižykite galimybių arba tikimybių medį.
\item Žodis ``ir'' dažnai reiškia sankirtą, žodis ``arba'' -- sąjungą.
\item Be grąžinimo tikimybės keičiasi, su grąžinimu -- paprastai išlieka tokios pačios.
\item Patikrinkite, ar įvykiai nesutaikomi arba nepriklausomi, prieš taikydami specialias formules.
\end{enumerate}

\subsection{Iliustruojantys pavyzdžiai}
\begin{enumerate}[label=\arabic*.]
\item Metamas sąžiningas kauliukas. Raskime tikimybę gauti lyginį skaičių.
Palankios baigtys: \(2,4,6\). Iš viso \(6\) baigtys.
\[
P=\frac36=\frac12.
\]
\textbf{Atsakymas:} \(\frac12\).

\item Dėžėje yra \(4\) balti ir \(3\) juodi rutuliai. Traukiami du be grąžinimo. Raskime tikimybę, kad abu balti.
\[
P=\frac47\cdot\frac36=\frac{12}{42}=\frac27.
\]
\textbf{Atsakymas:} \(\frac27\).

\item Tikimybės \(P(A)=0{,}6\), \(P(B)=0{,}5\), \(P(A\cap B)=0{,}2\). Raskime \(P(A\cup B)\).
\[
P(A\cup B)=0{,}6+0{,}5-0{,}2=0{,}9.
\]
\textbf{Atsakymas:} \(0{,}9\).
\end{enumerate}

\subsection{Dažnos klaidos}
\begin{itemize}
\item Klasikinė formulė taikoma nevienodai tikėtinoms baigtims.
\item Žodis ``arba'' suprantamas kaip visada nesutaikomi įvykiai, nors kartais įvykiai gali įvykti kartu.
\item Be grąžinimo antrame žingsnyje nepakeičiamas bendras baigčių skaičius.
\end{itemize}

\subsection{Mini praktika}
\begin{enumerate}[label=\alph*)]
\item Iš kortelių \(1,2,3,4,5\) atsitiktinai traukiama viena kortelė. Raskite tikimybę gauti pirminį skaičių.
\item Dėžėje \(2\) raudoni ir \(5\) mėlyni rutuliai. Traukiami du su grąžinimu. Raskite tikimybę, kad abu raudoni.
\item Jei \(A\) ir \(B\) nepriklausomi, \(P(A)=0{,}4\), \(P(B)=0{,}3\), raskite \(P(A\cap B)\).
\end{enumerate}

\subsection{Ryšys su kitomis temomis}
Tikimybiniai modeliai remiasi kombinatorika, aibėmis ir loginiais jungtukais. Jie paruošia atsitiktinius dydžius, binominį skirstinį ir statistinę išvadą.

\section{Atsitiktiniai dydžiai ir normalusis skirstinys}\label{sec:atsitiktiniai-dydziai-ir-normalusis-skirstinys}

\subsection{Apibrėžtys}
\begin{apibreztis}
Atsitiktinis dydis yra skaitinė bandymo rezultato funkcija, pavyzdžiui, herbų skaičius metant monetą arba laukimo laikas stotelėje.
\end{apibreztis}
\begin{apibreztis}
Normalusis skirstinys yra tolydus varpo formos skirstinys, žymimas \(X\sim N(\mu;\sigma^2)\), kur \(\mu\) yra vidurkis, o \(\sigma\) -- standartinis nuokrypis.
\end{apibreztis}
\begin{apibreztis}
Standartizavimas paverčia \(X\sim N(\mu;\sigma^2)\) standartiniu normaliuoju dydžiu:
\[
Z=\frac{X-\mu}{\sigma}.
\]
\end{apibreztis}

\subsection{Teoremos ir savybės}
\begin{savybe}
Normaliojo skirstinio kreivė simetriška apie \(\mu\). Didžiausias tankis yra ties vidurkiu.
\end{savybe}
\begin{savybe}
Apytikslė \(68\)-\(95\)-\(99{,}7\) taisyklė:
apie \(68\%\) reikšmių patenka į \([\mu-\sigma;\mu+\sigma]\), apie \(95\%\) -- į \([\mu-2\sigma;\mu+2\sigma]\), apie \(99{,}7\%\) -- į \([\mu-3\sigma;\mu+3\sigma]\).
\end{savybe}
\begin{savybe}
Tikimybės tolydžiajame skirstinyje atitinka plotus po tankio kreive, todėl \(P(X=a)=0\), o skaičiuojamos intervalų tikimybės.
\end{savybe}

\subsection{Taisyklės ir algoritmai}
\begin{enumerate}[label=\arabic*.]
\item Nustatykite \(\mu\) ir \(\sigma\).
\item Jei reikia, standartizuokite ribas: \(z=\frac{x-\mu}{\sigma}\).
\item Naudokite normalaus skirstinio lentelę, skaičiuoklę arba \(68\)-\(95\)-\(99{,}7\) taisyklę.
\item Atsakymą interpretuokite kaip tikimybę arba procentą.
\item Patikrinkite, ar gautas procentas logiškas pagal atstumą nuo vidurkio.
\end{enumerate}

\subsection{Iliustruojantys pavyzdžiai}
\begin{enumerate}[label=\arabic*.]
\item Testo rezultatai pasiskirstę normaliai, \(\mu=70\), \(\sigma=10\). Kokia apytikslė dalis rezultatų tarp \(60\) ir \(80\)?
Intervalas yra \(\mu\pm\sigma\), todėl pagal taisyklę apie \(68\%\).
\textbf{Atsakymas:} apie \(68\%\).

\item Tame pačiame modelyje standartizuokime rezultatą \(85\).
\[
z=\frac{85-70}{10}=1{,}5.
\]
\textbf{Atsakymas:} rezultatas yra \(1{,}5\) standartinio nuokrypio virš vidurkio.

\item Dydžio \(X\) skirstinio lentelė:
\[
\begin{array}{c|ccc}
x&0&1&2\\ \hline
P(X=x)&0{,}25&0{,}50&0{,}25
\end{array}
\]
Raskime viltį ir dispersiją.
\[
E(X)=0\cdot0{,}25+1\cdot0{,}50+2\cdot0{,}25=1.
\]
\[
D(X)=(0-1)^2\cdot0{,}25+(1-1)^2\cdot0{,}50+(2-1)^2\cdot0{,}25=0{,}5.
\]
\textbf{Atsakymas:} \(E(X)=1\), \(D(X)=0{,}5\).
\end{enumerate}

\subsection{Dažnos klaidos}
\begin{itemize}
\item \(\sigma\) painiojamas su \(\sigma^2\). Žymėjime \(N(\mu;\sigma^2)\) antras parametras yra dispersija.
\item Skaičiuojama taško tikimybė \(P(X=a)\) tolydžiajame skirstinyje kaip teigiamas plotas.
\item Standartizuojant sukeičiamas ženklas: turi būti \(x-\mu\), ne \(\mu-x\), jei naudojama įprasta \(Z\) formulė.
\end{itemize}

\subsection{Mini praktika}
\begin{enumerate}[label=\alph*)]
\item Jei \(X\sim N(100;15^2)\), standartizuokite \(x=130\).
\item Pagal \(68\)-\(95\)-\(99{,}7\) taisyklę į kokį intervalą patenka apie \(95\%\) reikšmių, kai \(\mu=50\), \(\sigma=4\)?
\item Kodėl normaliuoju skirstiniu aprašomam dydžiui \(P(X=50)=0\), bet \(P(49<X<51)>0\)?
\end{enumerate}

\subsection{Ryšys su kitomis temomis}
Atsitiktiniai dydžiai ir normalusis skirstinys sujungia tikimybę, funkcijų grafikus ir statistinę išvadą. Normalusis modelis padeda interpretuoti matavimo paklaidas, testų rezultatus ir imčių vidurkius.
